\chapter{Road-conditioned synthesis of complete journeys}
\label{ch:native-method}

\section{The conditional generation task}
The output is a complete timed journey in the form \texttt{time,id,x,y,type}. A request supplies a road network $M$, an origin and destination, a participant class $c$, and the vehicle's declared road permissions. A scheduled bus additionally supplies its ordered public stop anchors. Intermediate reference coordinates, realised arrival time and future speeds are unavailable. The system samples a permitted route $R$, a duration $T$, and a motion realization along that route. A new random draw can change the route, the motion, or both.

Let $R(s)$ interpolate a constructed route at arc distance $s\in[0,L]$, and let $p(t)$ denote normalized progress. The complete system factors its sampling distribution as
\begin{align}
 R&\sim p_{\rho}(\,\cdot\mid M,\mathrm{OD},c,\mathrm{stops}),\nonumber\\
 T&\sim p_{\mathrm{dur}}(\,\cdot\mid L,c),\nonumber\\
 h&\sim p_{\eta}(\,\cdot\mid R,T,c),\nonumber\\
 p&\sim p_{\theta}(\,\cdot\mid R,T,c,h),\qquad
 X_i=R\bigl(Lp(t_i)\bigr).
 \label{eq:native-factorization}
\end{align}
The optional condition $h$ represents the learned spatial and speed views described in Chapter~\ref{ch:native-feedback}; the baseline omits it. Its distribution is fitted using training trajectories and sampled for new requests. At synthesis, the sampled code conditions temporal diffusion directly, producing route progress for coordinate decoding.

The factorization defines a modular generative model in which route selection establishes spatial support, the duration prior determines the journey's temporal scale, and diffusion models its progress profile. Road-coordinate decoding combines these outputs into a complete timestamped trajectory. This decomposition also supports controlled experiments: matching route and duration draws isolates changes in the generated motion profile.

\fig{g40-framework.pdf}{1}{Implemented training and out-of-sample synthesis. Observed training trajectories produce the two image views and generative conditions; new synthesis samples conditions from a training-only prior. Route choice, duration, motion diffusion and coordinate decoding have distinct responsibilities. The learned condition enters every denoising block, coupling representation learning to complete-journey generation.}{g40-framework}

Training uses observed route--motion pairs, whereas test-time routes are constructed from the road network and request attributes. This distinction makes the learned temporal model usable for new journeys within the available road catalog~\cite{difftraj,diffrn}.

\section{Learning route preferences within a permitted catalog}
Road legality alone leaves a choice between several valid ways of travelling from an origin to a destination. The route module first constructs a finite catalog, then assigns learned probabilities to its members. Construction uses the directed road graph and the declared SUMO vehicle permissions. Road access follows the declared vehicle permissions, including those assigned to taxi/ride-hailing vehicles. For buses, successive public stop anchors divide the route into ordered subproblems.

The constructor makes at most twelve deterministic proposals using shortest-path search with penalties that encourage alternatives. It retains at most four geometrically supported paths and discards candidates longer than 1.5 times the first valid route. Actual lane connections, internal junction geometry and endpoint positions are checked during construction. The resulting catalog defines the finite spatial support over which the learned policy assigns probabilities.

The policy is fitted by behavioural cloning of permitted branching decisions in training trips. At an edge $e$ with next public goal $g$, it considers directed successors from which $g$ is reachable. Each successor $a$ has twenty static features: current and successor coordinates, goal coordinates and relative goal vectors, edge length, speed limit, lane count, turn and goal alignment, remaining shortest-path distances, excess path cost and a goal-arrival indicator. Their physical scales are fixed. These features describe network geometry and the current routing goal.

A multilayer perceptron with widths $20\rightarrow128\rightarrow64\rightarrow5$ produces a utility for each class. Selecting the requested class output and masking unsupported successors gives
\begin{equation}
 \pi_\rho(a\mid e,g,c,M)=
 \frac{\exp u_{\rho,c}(f(e,a,g,M))}
 {\sum_{b\in\mathcal A(e,g,M)}\exp u_{\rho,c}(f(e,b,g,M))}.
 \label{eq:route-policy}
\end{equation}
Cross-entropy fitting samples a class, source vehicle, leg and branching decision uniformly at each level. The supported training subset contains 23,696 legs and 546,363 branching decisions; the corresponding development subset contains 5,748 legs and 132,341 decisions. Classification training uses transitions with multiple reachable successors. Two fixed neural fits are retained, alongside linear and shortest-cost controls in the route study.

For a catalog route $r$, the score is the sum of its branching-choice log probabilities, with the goal updated at ordered stop anchors. Normalizing these scores over the available catalog yields
\begin{equation}
 p_\rho(r\mid M,\mathrm{OD},c,\mathrm{stops})=
 \frac{\exp\!\left\{\sum_{j\in\mathcal B(r)}
 \log\pi_\rho(e_{j+1}\mid e_j,g_j,c,M)\right\}}
 {\sum_{r'\in\mathcal C}\exp\!\left\{\sum_{j\in\mathcal B(r')}
 \log\pi_\rho(e_{j+1}\mid e_j,g_j,c,M)\right\}}.
 \label{eq:route-catalog}
\end{equation}
The branching log probabilities are summed directly and normalized over the catalog, yielding a behavioral policy informed by recursive route-choice modeling~\cite{routechoice}. An unsupported policy evaluation uses a recorded uniform fallback over the valid catalog.

Route quality is consequently assessed through route-use distributions, coverage and alternative-path diversity, together with the generated motion. Distributional evaluation accommodates multiple permitted alternatives between the same endpoints. In the main conditioning comparison, the route policy and its random keys are fixed across all model variants. This pairing isolates changes in movement conditional on the selected route.

\section{Representing motion on physical time}
For a training trajectory, let $\ell_i=\|X_i-X_{i-1}\|_2$, $s_i=\sum_{j\leq i}\ell_j$ and $p_i=s_i/L$, where $L=s_n>0$. Repeated positions give repeated progress values and therefore retain stopped intervals. All nonduplicate spatial vertices are stored for exact piecewise-linear reconstruction; a separate set of 129 equally spaced arc positions provides the network's route description. The network thus receives a compact route summary while decoding retains the full polyline.

Two temporal representations are compared. The relative-time version uses 257 points uniformly distributed over $[0,T]$. The physical-clock version uses $0,0.25,0.50,\ldots,T$, appending the exact final time when it is not a multiple of 0.25 seconds. Its number of intervals varies with duration. The distinction matters because a fixed number of controls gives a long trip much coarser physical resolution than a short one. The temporal-resolution experiment measures how this representation choice affects generated movement.

With $\Delta t_i=t_i-t_{i-1}$, define the dimensionless relative speed and its transformed target as
\begin{equation}
 r_i=\frac{p_i-p_{i-1}}{\Delta t_i/T},\qquad
 u_i=\log(1+r_i),\qquad z_{0,i}=\frac{u_i-\mu}{\sigma}.
 \label{eq:native-target}
\end{equation}
Thus $r_i=1$ is constant progress at the journey's average rate, and $r_i=0$ represents a stop. The logarithm moderates the upper tail without removing zero motion. Training-only $\mu$ and $\sigma$ assign equal mass to each trajectory and weight its intervals by $\Delta t_i/T$. Both clock versions use this target definition, with their own fitted normalization. Together, the two representations form a matched comparison on the current network.

\section{Class-Conditioned Travel Duration Prior}
Within each class, a ridge regression relates log travel time to centred log route length. For training journey $j$ of class $c$,
\begin{equation}
 \log T_j=a_c+b_c(\log L_j-\bar\ell_c)+e_j,
 \qquad
 T^*=\exp\{a_c+b_c(\log L-\bar\ell_c)+e^*\},
 \label{eq:native-duration}
\end{equation}
where $e^*$ is sampled uniformly from that class's training residuals. The slope has ridge penalty $10^{-3}$; the intercept is unpenalized. Residual sampling retains variation beyond the conditional mean without imposing a Gaussian residual law. A numerical lower bound gives $T=\max(T^*,0.25,L/55)$, with metres and seconds as the physical units. This lower bound provides a numerical guard for duration sampling.

In the temporal-resolution pilot, both diffusion versions and the empirical comparator share the same prior and request-keyed draws. The main conditioning experiment instead holds the expanded training-side prior fixed across all eight model variants, with matched route and duration draws. Consequently their unselected duration distributions are paired by construction. The shared prior makes the motion mechanisms the source of differences in generated speed variation. Terminal feedback may select different durations from the shared pool; that is a selection effect and is reported as such.

\section{A masked temporal denoiser}
The base denoiser is a one-dimensional U-Net with channel widths 64, 96 and 128 and 891,969 parameters. The conditioning framework includes an additional 24,832-parameter condition adapter; the B variant bypasses it. The separately trained image encoders are accounted for in Chapter~\ref{ch:native-feedback}. Two downsampling stages, skip connections and nearest-neighbour upsampling connect the resolutions. Residual convolutional blocks use dilations 1 and 2, with dilations 1, 2, 4 and 8 in the bottleneck. Each block also receives a masked global feature summary, allowing information to pass beyond its local receptive field.

The conditioning vector adds a sinusoidal diffusion-step embedding, a learned class embedding, a convolutional route embedding and a five-component context embedding. The context contains normalized starting coordinates, log route length, log sampled duration and log mean travel speed. Route coordinates are expressed relative to the start and scaled by 1,500 metres. Three channels describe each interval's physical midpoint, relative midpoint and actual width. These allow the network to distinguish physical time from relative trip position. The denoiser learns stopping behavior from the motion targets, using route geometry and the sampled duration as context.

Sequences are padded to a multiple of four. The mask is reapplied through the network, pooling divides by the number of valid entries, and normalization acts across channels at each position rather than across padded time. The mask restricts computation and loss to the observed temporal support. Batches are sampled before length sorting; sorting and microbatching manage memory without removing long journeys.

Following noise-prediction diffusion~\cite{ddpm}, a uniformly sampled noise index $k$ corrupts a target by
\begin{equation}
 z_k=\sqrt{\bar\alpha_k}z_0+\sqrt{1-\bar\alpha_k}\,\epsilon,
 \qquad \epsilon\sim\mathcal N(0,I),
 \label{eq:native-forward}
\end{equation}
on valid intervals. The 1,000-step schedule increases $\beta_k$ linearly from $10^{-4}$ to $0.02$, with $\bar\alpha_k=\prod_{j\leq k}(1-\beta_j)$. For an effective batch of $B$ trajectories, the loss is
\begin{equation}
 \mathcal L=\frac{1}{B}\sum_{b=1}^{B}\sum_i
 \frac{\Delta t_{b,i}}{T_b}
 \left[\epsilon_{b,i}-\epsilon_\theta(z_{b,k},k,R_b,c_b,T_b,h_b)_i\right]^2.
 \label{eq:native-loss}
\end{equation}
Padding has zero weight. Physical-time weighting within a trajectory prevents the fractional final interval from receiving a full interval's influence; averaging trajectories prevents long trips from dominating merely through their larger point counts.

\section{Sampling and coordinate decoding}
Generation uses 100 deterministic DDIM updates~\cite{ddim} between equally spaced indices of the training schedule. At a current index $k$, the predicted clean target and next state at $j<k$ are
\begin{align}
 \hat z_0&=(z_k-\sqrt{1-\bar\alpha_k}\,\hat\epsilon)/\sqrt{\bar\alpha_k},\\
 z_j&=\sqrt{\bar\alpha_j}\,\hat z_0+\sqrt{1-\bar\alpha_j}\,\hat\epsilon.
 \label{eq:native-ddim}
\end{align}
The final update uses $\bar\alpha_j=1$. Different initial Gaussian draws produce candidates; paired request and candidate identifiers reproduce duration draws across methods. Each temporal grid has its own noise-array dimensions.

After undoing standardization, the decoder clips $u_i$ to $[0,8]$ and constructs nonnegative progress mass:
\begin{equation}
 m_i=\left[\exp\{\operatorname{clip}(\mu+\sigma\hat z_{0,i},0,8)\}-1\right]
       \frac{\Delta t_i}{T},\qquad
 \hat p_i=\frac{\sum_{j\leq i}m_j}{\sum_jm_j}.
 \label{eq:native-decoder}
\end{equation}
The initial progress is zero and the final progress is one. If total mass is below $10^{-12}$, the decoder falls back to constant progress in physical time. Decoder diagnostics record negative-value truncation, upper clipping and constant-progress fallback. The normalization couples the whole journey: a local mass change can alter progress elsewhere. Endpoint completion follows from normalization, while the learned mass profile determines acceleration and stop timing.

Coordinates are interpolated along all retained vertices of the constructed map path. The motion model generates longitudinal progress along the selected route. For the relative-time model, progress is first interpolated onto the same quarter-second output clock; the exact fractional endpoint remains. Thus both methods deliver comparable timestamped coordinates. Permitted lanes, directed connections and source stop anchors are handled by the route constructor. This geometric constraint is shared with the empirical baseline and all GAF selectors.

The baseline generates complete journeys using class, route and duration alone. The GAF-conditioned variants retain the same target, physical-time grid, sampler and decoder, while adding a sampled condition at every denoising step. Chapter~\ref{ch:native-feedback} describes how that condition is learned and generated without future trajectory access. Earlier terminal-selection experiments remain separate controls on this synthesis system.
